\documentclass[11pt,a4paper]{article}
% Common preamble for the hint sheets (\input from each file)
\usepackage[utf8]{inputenc}
\usepackage[T1]{fontenc}
\usepackage{lmodern}
\usepackage{amsmath,amssymb,amsthm,mathtools}
\usepackage[a4paper,margin=2.3cm]{geometry}
\usepackage{enumitem}
\usepackage{xcolor}
\usepackage[hidelinks]{hyperref}
\renewcommand{\leq}{\leqslant}
\renewcommand{\le}{\leqslant}
\renewcommand{\geq}{\geqslant}
\renewcommand{\ge}{\geqslant}
\setlength{\parindent}{0pt}
\setlength{\parskip}{0.35em}
% \hintstitle{Sheet title}{Class work / Homework}
\newcommand{\hintstitle}[2]{\begin{center}{\LARGE\bfseries Hints: #1\par}\vspace{0.5em}{\large #2\par}\end{center}
\noindent\textit{The hints to each problem are cumulative. Read one, try again, and open the next only if you are still stuck.}\par\medskip\hrule\medskip}
\newcommand{\problem}[1]{\par\bigskip{\large\bfseries Problem #1}\par\nopagebreak\smallskip}
\newcommand{\subpart}[1]{\par\smallskip\textbf{(#1)}\par\nopagebreak}
\newcommand{\hint}[1]{\par\smallskip\noindent\textcolor{black!70}{\textbf{Hint #1.}}\enspace}
\newcommand{\note}[1]{\par\smallskip\noindent\textit{Note.} #1\par}
\providecommand{\lt}{<}
\providecommand{\gt}{>}
\providecommand{\tightlist}{\setlength{\itemsep}{0pt}\setlength{\parskip}{0pt}}
\emergencystretch=3em

\begin{document}
\hintstitle{Abstract algebra -- 3. Permutations}{Class work}

\problem{1}
\hint{1} The cycles of $\sigma$ are forced: they are the orbits of $\langle\sigma\rangle$, and on each orbit $\sigma$ acts as one cycle. Disjoint cycles commute, so $\sigma^k=\mathrm{id}$ iff $c^k=\mathrm{id}$ for every cycle $c$, i.e.\ iff every cycle length divides $k$.
\hint{2} Conjugacy: $\tau(i_1\dots i_k)\tau^{-1}=(\tau(i_1)\dots\tau(i_k))$. So conjugate permutations have the same cycle type. Conversely, two permutations of the same type are conjugate by the $\tau$ that sends one cycle to the other, cycle by cycle.
\hint{3} The largest order in $S_n$ is $g(n)=\max\operatorname{lcm}(\lambda_1,\dots,\lambda_k)$ over partitions of $n$. You may assume the parts are prime powers with distinct primes. For $n=10$ the maximum is $2\cdot3\cdot5=30$. For large $n$, $\ln g(n)\sim\sqrt{n\ln n}$ (Landau): use primes up to about $\sqrt{n\ln n}$.

\problem{2}
\hint{1} Show that multiplying $\sigma$ by a transposition $(ab)$ changes the number of cycles by exactly $\pm1$: it merges the two cycles if $a,b$ lie in different cycles and splits the cycle if they lie in the same one. It also changes the parity of $\operatorname{inv}$: an adjacent swap changes $\operatorname{inv}$ by $\pm1$, and any transposition is a product of an odd number of adjacent swaps.
\hint{2} Write $\sigma$ as a product of transpositions starting from $\mathrm{id}$. Both $(-1)^{\operatorname{inv}}$ and $(-1)^{n-c}$ flip at every step, so they agree. The second is clearly multiplicative: a product of $k$ transpositions has sign $(-1)^k$, so the sign is a homomorphism.
\hint{3} The reversal $k\mapsto n+1-k$ has all $\binom n2$ pairs as inversions, so its sign is $(-1)^{n(n-1)/2}$. This is $+1$ iff $n\equiv0,1\pmod4$.

\problem{3}
\hint{1} $\pi$ fixes $0$, and $\mathbb F_p^*$ is cyclic of order $n=p-1$ (generator $g$). Writing $x=g^k$, the permutation becomes $k\mapsto3k$ on $\mathbb Z/n$, which is a bijection because $3\nmid n$.
\hint{2} Use the Chinese remainder theorem: $\mathbb Z/n\cong\mathbb Z/2^a\times\mathbb Z/m$ with $m$ odd, and multiplication by $3$ acts on each factor. For $\sigma\times\tau$ acting on $X\times Y$, $\operatorname{sgn}=\operatorname{sgn}(\sigma)^{|Y|}\operatorname{sgn}(\tau)^{|X|}$. So the sign is that of $k\mapsto3k$ on $\mathbb Z/2^a$.
\hint{3} If $a=1$ (that is, $p\equiv3\pmod4$), this map is the identity on $\mathbb Z/2$: even. If $a\ge2$, count its cycles using $\operatorname{ord}_{2^j}(3)$ ($=1,1,2$ for $j=0,1,2$ and $2^{j-2}$ for $j\ge3$): there are $3+2(a-2)$ cycles, so $2^a-c$ is odd and the map is odd.

\problem{4}
\hint{1} Write $\frac1{\nu+1}=\int_0^1x^{\nu}\,dx$. Then the sum is $\int_0^1\sum_\pi\sigma(\pi)x^{\nu(\pi)}\,dx$.
\hint{2} $\sum_\pi\sigma(\pi)x^{\nu(\pi)}$ is the determinant of the matrix with $x$ on the diagonal and $1$ elsewhere, i.e.\ $\det((x-1)I+J)=(x-1)^{n-1}(x-1+n)$ (the eigenvalues of $J$ are $n,0,\dots,0$).
\hint{3} Integrate $(x-1)^n+n(x-1)^{n-1}$ from $0$ to $1$.

\problem{5}
\hint{1} Use linearity of expectation with indicators: $k$ is fixed with probability $1/n$. For cycles, build $\sigma$ by the ``closing a cycle'' construction (or count): the $k$-th step closes a cycle with probability $1/k$, so the expectation is $H_n$.
\hint{2} Derangements by inclusion--exclusion: $D_n/n!=\sum_{k=0}^n(-1)^k/k!\to e^{-1}$.
\hint{3} Local maxima: an end position is a local maximum with probability $1/2$, an interior one with probability $1/3$ (the largest of three random values). For $n\ge2$ the answer is $2\cdot\frac12+\frac{n-2}3=\frac{n+1}3$.

\end{document}
